\chapter{BF Theory}
\label{ch:tqft-bf-theory}

\ConventionsEuclidean

BF theory is the first cohomological gauge theory developed in this volume.
It is simple enough to display the fields, equations, gauge symmetries,
observables, and BV reducibility explicitly, while still showing why
topological gauge theory is more than a number assigned to a closed manifold.

\section{Classical Fields and Action}

Let \(M\) be an oriented smooth \(D\)-manifold, and let \(P\to M\) be a
principal \(G\)-bundle for a compact Lie group \(G\).  Fix an invariant
pairing \(\langle\cdot,\cdot\rangle\) between \(\mathfrak g^*\) and
\(\mathfrak g\).  The classical fields of nonabelian BF theory are
\[
  A\in\Omega^1(M,\operatorname{ad}P),
  \qquad
  B\in\Omega^{D-2}(M,\operatorname{ad}P^*).
\]
The curvature is \(F_A=dA+\frac12[A,A]\).  The action is
\[
  S_{\rm BF}[A,B]
  =
  2\pi\ii\int_M \langle B\wedge F_A\rangle ,
\]
with the coefficient chosen so that integral-period constraints in abelian
compact examples can be imposed without changing the local formula.

Varying \(B\) gives
\[
  F_A=0.
\]
Varying \(A\) gives
\[
  d_A B=0.
\]
Thus the stationary points are flat connections equipped with covariantly
closed \(B\)-fields, modulo the gauge symmetries described below.

\section{Gauge Symmetries and Reducibility}

Ordinary gauge transformations act by
\[
  A\mapsto gAg^{-1}-dg\,g^{-1},
  \qquad
  B\mapsto gBg^{-1}
\]
after identifying \(\mathfrak g^*\) with \(\mathfrak g\) by an invariant
form.  Infinitesimally, with \(c\in\Omega^0(M,\operatorname{ad}P)\),
\[
  \delta_c A=d_Ac,
  \qquad
  \delta_c B=[B,c].
\]
There is also a shift symmetry
\[
  \delta_\eta A=0,
  \qquad
  \delta_\eta B=d_A\eta,
  \qquad
  \eta\in\Omega^{D-3}(M,\operatorname{ad}P^*).
\]
On the locus \(F_A=0\), this shift symmetry is reducible because
\(\eta=d_A\eta_1\) gives \(\delta_\eta B=d_A^2\eta_1=0\).  The reducibility
continues through a chain
\[
  \Omega^{D-4}\to\Omega^{D-3}\to\Omega^{D-2}.
\]
Therefore a BV description is not decoration; it is the natural field space
for ghosts, ghosts-for-ghosts, antifields, and the master equation.

\section{Metric Independence}

The displayed action uses only the orientation of \(M\), the exterior
derivative, wedge product, bundle data, and invariant pairing.  It contains no
metric.  If the gauge fixing and regularization can be chosen so that their
metric dependence is BV-exact and no anomaly remains, correlation functions
of BV-closed observables are metric independent.  The verification has two
parts: the classical action has no metric dependence, and the regulated
quantum measure plus gauge fixing has vanishing metric anomaly in the sector
being used.

\section{Observables}

For a closed loop \(\gamma\subset M\) and a representation \(R\) of \(G\), the
Wilson loop
\[
  W_R(\gamma)=\operatorname{Tr}_R\,\operatorname{Hol}_\gamma(A)
\]
is gauge invariant.  In abelian BF theory, if \(G=U(1)\), the field \(B\) is a
\((D-2)\)-form gauge field and a closed \((D-2)\)-cycle \(\Sigma\) defines
\[
  V_n(\Sigma)=\exp\!\left(2\pi\ii n\int_\Sigma B\right),
  \qquad n\in\mathbb Z,
\]
when \(B\) has the corresponding compact-period normalization.  The path
integral pairs electric Wilson loops and magnetic \(B\)-surface operators by
intersection or linking data.  The precise formula depends on the global
normalization of the compact fields and must be derived from the finite
abelian gauge group or differential-cohomology model being used.

\begin{example}[Finite abelian BF theory]
\label{ex:finite-abelian-bf}
For a finite abelian group \(A\), the fully finite theory can be defined by
summing over flat \(A\)-bundles.  On a closed \(D\)-manifold \(M\), the
partition function is a normalized count of maps from the fundamental
groupoid of \(M\) to \(A\), with gauge automorphisms divided out.  This
definition is a genuine finite sum and supplies a rigorous model for the
topological sector that continuum \(U(1)\) BF theory approximates after
choosing compact periods.
\end{example}

\begin{openproblem}[Continuum BF quantization with boundaries]
\label{op:continuum-bf-boundaries}
Construct nonabelian continuum BF theory on manifolds with boundary as a
functorial BV theory.  The construction must specify the boundary phase
space, residual fields, gauge fixing, anomaly line, gluing pairing, and the
relationship between Wilson, surface, and boundary observables.
\end{openproblem}
